\chapter{ENGINEERING CONSIDERATIONS, CHALLENGES
AND REMEDIES}
\onehalfspacing
\label{chap:engineering-considerations}
\vspace{-3em}

This chapter gives a clear and detailed look at the engineering ideas, ethical duties and technical challenges we faced while building our \textbf{LLM-Driven Clinical QnA Classification} system. This system represents a new approach to screening for Borderline Personality Disorder (BPD) using natural language processing. Here, we'll explain how we tackled these points through carefully planned engineering steps, focusing especially on keeping \textbf{participant data private}, using AI responsibly in a sensitive medical setting and managing the difficulties that came with creating such a unique screening tool. Understanding these aspects is crucial for appreciating the depth and thoughtfulness behind our project.

\section{Ethical Considerations}
Our system was designed to sort answers to medical questions for BPD screening which naturally handles private health information. This meant ethics were at the forefront of our work. We went to great lengths to ensure all data collection strictly followed ethical rules. For every piece of information gathered through Google Forms, we \textbf{got clear and informed consent from participants}. This involved making sure they understood what data was being collected, how it would be used and that their participation was completely voluntary. Crucially, all collected data was immediately handled and used in a way that made it \textbf{anonymous}, meaning no personal details could ever be linked back to an individual, thereby keeping their privacy safe.

Ethical development also meant making sure our system wouldn't produce misleading results, continue any unfair leanings, or accidentally lead to wrong diagnoses or negative labeling of individuals. We took extreme care during the system's development and testing to ensure its classification results consistently matched what is already known and accepted in clinical practice. It's vital to stress that this system is developed as a \textbf{screening tool} only, meant to help identify potential indicators and is emphatically not a diagnostic instrument. It should never replace a professional medical assessment. Throughout our entire research process, we strictly followed honest academic rules, being completely open about our methods and findings and always taking full responsibility for our work. Our highest priority was always the \textbf{participant well-being} and maintaining the absolute safety and integrity of their data.

\section{Data Security Consideration}
As clinical data is so sensitive and private, we made data security and confidentiality our absolute top priority from day one. When we collected any raw data, it was immediately \textbf{made unidentifiable and had all personal details removed}. This process of de-identification ensures that even if the data were somehow accessed, it couldn't be traced back to any individual.

All the later steps in our project like preparing the data for use, training our machine learning model and finally running it to get results were done \textbf{on our own computers in a safe, controlled environment}. This 'local' approach was a deliberate choice to avoid the risks associated with transmitting or storing sensitive data on outside online services or cloud platforms. We did not use any external APIs or cloud tools for processing or storing any raw, identifiable \textbf{participant data} at any point. Furthermore, the entire dataset was put together and organized offline to prevent any accidental exposure to the internet or other external threats. This rigorous approach made absolutely sure that all \textbf{participant information} stayed confidential and protected throughout the entire research process, from collection to analysis.


\section{Complex Engineering Problem}
Our thesis project tackles many complex engineering issues that align directly with the detailed criteria outlined in the OBE (Outcome-Based Education) guidelines. There are seven specific criteria for complex problem-solving in total and we successfully addressed the following four in our study.

\subsection{P1: Deep Knowledge Needed}
Building this system required us to bring together several advanced and specialized fields of knowledge. We had to combine sophisticated natural language processing methods, techniques for fine-tuning transformer models (which are a type of deep learning model), careful ways to prepare clinical text for analysis and reliable classification techniques. This wasn't just about knowing theories; it truly \textbf{needed a deep understanding} of these areas in practice. For instance, understanding \textbf{clinical NLP} meant knowing how to process medical jargon, abbreviations and the subtle ways people express themselves in healthcare contexts. Our knowledge of \textbf{deep learning designs} allowed us to choose and adapt the right transformer architecture for our specific QnA task. Finally, knowing \textbf{how psychiatric assessments work} guided us in structuring our data and interpreting our model's outputs in a clinically meaningful way.

\subsection{P4: Understanding the Challenges}
Working with medical question-and-answer data presented significant challenges because it was often messy, subtle and heavily dependent on context for accurate psychiatric screening. This meant we needed a deep familiarity with both new machine learning methods and the \textbf{specific difficulties posed by medical language and the symptoms of BPD}. For example, data might be "uneven" where some categories have far more examples than others which can make models struggle to learn. We had to develop strategies to deal with this, like special sampling techniques. We also had to recognize "small language clues" tiny words or phrases that might indicate a symptom but are easily missed by general NLP tools. Furthermore, we had to be constantly aware of the risk of our model developing "unfair leanings" or biases and actively work to prevent this, especially given the sensitive nature of BPD screening.

\subsection{P6: How Many People Were Involved \& Their Different Needs}
Creating a medical screening tool isn't just a technical task; it meant actively thinking about what different groups of people needed and expected. These stakeholders included \textbf{participants} who provided the data, \textbf{doctors} who might use such a tool and \textbf{researchers} who needed clear, verifiable results. Often, these groups had different needs that could "clash." For example, doctors might want a model that's very easy to understand (interpretability), even if it's slightly less accurate, while researchers might prioritize the highest possible classification accuracy. Similarly, there was a tension between wanting the tool to work for a wide range of situations versus needing very precise results for a specific medical condition. Finding a balanced technical design that could meet these various needs related to how useful the tool was, how safe it was and its ethical impact was a continuous challenge.

\subsection{P7: How Parts Depend on Each Other}
The clinical QnA classification system we built was made of several parts that were \textbf{closely connected and relied heavily on each other}. These included: the process for \textbf{secure data collection and anonymization}, the special NLP feature extraction steps, the training and fine-tuning of the transformer model itself and finally, the classification part that made the predictions. Because these components were so interdependent, if one part didn't work perfectly, it directly affected how well the whole system performed and how dependable it was. For instance, if the data collection wasn't meticulous, or if the anonymization step missed something, it could compromise privacy. If the NLP feature extraction wasn't robust, the model training would suffer, leading to inaccurate classifications. It was like a chain reaction where each link had to be strong for the entire system to succeed.

\clearpage

\section{Engineering Activities}
There are five main types of engineering activities. In our study, we actively carried out the following three.

\subsection{Types of Resources Used}
Our system utilized a specific but varied set of resources, chosen to be efficient and effective. This included substantial \textbf{computing power for training transformer models}, often relying on specialized hardware like a powerful GPU setup. We also leveraged publicly available NLP tools and libraries, such as those from the Hugging Face ecosystem which provided a strong foundation. Beyond off-the-shelf tools, we developed many \textbf{custom scripts for data anonymization and preprocessing} to specifically handle the unique characteristics of our clinical QnA data. All of this was integrated using standard machine learning frameworks, ensuring a streamlined development process.

\subsection{Innovation}
The core innovation of our work lies in the \textbf{new way we applied transformer-driven QnA classification as a dedicated NLP-based method for BPD screening}. While transformers are known for their power in NLP, their specific application as a focused screening tool for a complex mental health condition like BPD, particularly through QnA classification, represents a unique contribution. This involved making advanced NLP models suitable for a specific medical setting and creating a classification process made specifically to understand and interpret the unique features of medical conversation data related to BPD symptoms. This tailored approach allowed us to address the nuances of clinical language in a way general NLP models might miss.

\subsection{Familiarity}
Even though using these components together for BPD screening was a new and unique project, our work strongly relied on well-known and established engineering methods. We built upon solid foundations in \textbf{natural language processing}, the principles of \textbf{building deep learning models}, the use of \textbf{strong and reliable classification algorithms} and the best practices of \textbf{designing software in separate, manageable parts}. This reliance on established practices, even when venturing into novel applications, gave our work a solid, dependable base and allowed us to tackle challenges with confidence.

Ultimately, we engineered this system with clear, practical solutions to problems that were both very complex and carried high responsibility. All of this was done within a sensitive clinical area, truly showcasing an innovative and robust approach to BPD screening.

